\documentclass[11pt,a4paper]{article}
\usepackage[utf8]{inputenc}
\usepackage[T1]{fontenc}
\usepackage[french]{babel}
\usepackage[margin=2.2cm]{geometry}
\usepackage{amsmath,amssymb}
\usepackage{booktabs}
\usepackage{hyperref}
\usepackage{enumitem}

\title{Experiments v2 : contraintes structurelles pour ameliorer le taux de planarite}
\author{Stage M2 IMD --- Generateur non-benzenoides}
\date{Mai 2026}

\begin{document}
\maketitle

\section{Motivation}

L'analyse statistique de \texttt{db\_v2} (8 configurations CSP existantes,
toutes valeurs de $h$) a fait apparaitre deux regularites empiriques
robustes sur le statut de planarite des molecules generees :

\begin{enumerate}[leftmargin=*,itemsep=2pt]
  \item \textbf{Symetrie 5/7.} Les molecules dont le nombre de pentagones
        et d'heptagones sont proches ($|n_5 - n_7| \le 1$) atteignent un
        taux de planarite de 25--43~\%, contre 0--7~\% quand l'ecart
        depasse 1.
  \item \textbf{Pentagones en bord.} Le taux de planarite chute de
        59~\% (zero pentagone en bord) a 2~\% (cinq ou plus pentagones
        en bord du squelette).
\end{enumerate}

L'objectif d'\texttt{experiments\_v2} est de \emph{traduire ces lois
empiriques en contraintes CSP a priori}, pour generer directement des
molecules de meilleure qualite plutot que de filtrer a posteriori.

\section{Formalisation}

Soit $G = (V,E)$ le squelette hexagonal d'une molecule cible, ou chaque
sommet $v \in V$ represente un hexagone. La variable de decision
$x_v \in \{5,6,7\}$ donne la taille du cycle assignee a $v$. Un sommet
est qualifie \emph{de bord} (\texttt{boundary}) si son degre dans le
graphe dual est strictement inferieur a 6 ; sinon il est interieur.
Notons $B \subset V$ l'ensemble des sommets de bord.

Les nouvelles contraintes ajoutees sont parametrees par des bornes
entieres positives :

\begin{align*}
  \text{(C-SYM)} \quad & \left| \textstyle\sum_v \mathbf{1}[x_v=5] - \sum_v \mathbf{1}[x_v=7] \right| \le K_{\text{sym}} \\
  \text{(C-PB)}  \quad & \sum_{v \in B} \mathbf{1}[x_v=5] \le K_{\text{pb}} \\
  \text{(C-HB)}  \quad & \sum_{v \in B} \mathbf{1}[x_v=7] \le K_{\text{hb}}
\end{align*}

Ces trois contraintes sont injectees dans le modele PyCSP3 existant
(\texttt{csp\_model.py}) en plus des contraintes geometriques deja
presentes (frequences globales, adjacences 5--7 interdites, etc.). Le
solveur ACE reste inchange.

\section{Pipeline et infrastructure}

L'arborescence \texttt{experiments\_v2/} duplique \texttt{experiments/}
pour eviter toute regression sur le pipeline historique. Sept presets
sont definis dans \texttt{configs.py} :
\texttt{baseline\_v2}, \texttt{sym1} ($K_{\text{sym}}{=}1$),
\texttt{pb2} ($K_{\text{pb}}{=}2$), \texttt{sym1\_pb2},
\texttt{all\_strict} ($K_{\text{sym}}{=}0,\ K_{\text{pb}}{=}2,\ K_{\text{hb}}{=}3$),
\texttt{pb1}, \texttt{pb0}.

\paragraph{Mode DB-only.} Pour eviter les millions de fichiers \texttt{.xyz}
sur le NFS, chaque worker ecrit directement une mini base SQLite
(\texttt{worker\_dbs/<h>\_\_<config>\_\_<mol>.db}) contenant les
geometries en \texttt{BLOB gzip}. Un script \texttt{finalize.py}
fusionne ces bases en une seule \texttt{db\_v3.db} via
\texttt{ATTACH DATABASE} et \texttt{INSERT OR REPLACE} (idempotent). La
table \texttt{configs} est reconstruite globalement en post-merge.

\paragraph{Deploiement cluster.} Le pipeline reutilise l'infrastructure
SSH existante (\texttt{experiments/cluster/dispatcher.py}) en pointant
sur \texttt{experiments\_v2/cluster/worker.py}. Les 16 machines du
cluster COALA executent jusqu'a 20 jobs en parallele chacune. Toute
l'orchestration (claims atomiques sur NFS, timeouts, recuperation)
reste celle deja eprouvee.

\section{Resultats}

Les experimentations ont porte sur les graphes $h \in \{6,7,8,9\}$ de
\texttt{benzdb}, avec les flags \texttt{no-freeze} et \texttt{no-table}
actives (notre meilleur regime CSP de base). Le tableau 1 resume le
taux de planarite et le volume de solutions planaires obtenu pour les
trois contraintes \texttt{pb} testees sur tous les $h$.

\begin{table}[h]
\centering
\small
\begin{tabular}{lrrrrrr}
\toprule
$h$ & \multicolumn{2}{c}{\texttt{pb2}} & \multicolumn{2}{c}{\texttt{pb1}} & \multicolumn{2}{c}{\texttt{pb0}} \\
\cmidrule(lr){2-3}\cmidrule(lr){4-5}\cmidrule(lr){6-7}
    & \%plan & \#plans & \%plan & \#plans & \%plan & \#plans \\
\midrule
h6 & 68.5 & 2\,095   & \textbf{80.5} & 645    & --- (UNSAT) & 0 \\
h7 & 53.6 & 10\,443  & \textbf{68.9} & 2\,418 & --- (UNSAT) & 0 \\
h8 & 44.6 & 52\,799  & \textbf{61.7} & 9\,391 & --- (UNSAT) & 0 \\
h9 & 23.0 & 305\,605 & \textbf{31.3} & 40\,852 & --- (UNSAT) & 0 \\
\bottomrule
\end{tabular}
\caption{Taux de planarite et volume de solutions planaires par
configuration et par $h$ (no-freeze, no-table).}
\end{table}

Pour $h = 7$ specifiquement, la comparaison croisee des 7 presets a
egalement ete realisee :

\begin{table}[h]
\centering
\small
\begin{tabular}{lrr}
\toprule
Configuration & \#sols total & \%plan \\
\midrule
\texttt{baseline\_v2}                       & 29\,663 & 50.5 \\
\texttt{sym1}                                & 29\,663 & 50.5 \\
\texttt{pb2} / \texttt{sym1\_pb2} / \texttt{all\_strict} & 19\,478 & 53.6 \\
\texttt{pb1}                                 & 3\,510  & \textbf{68.9} \\
\texttt{pb0}                                 & 0       & UNSAT \\
\bottomrule
\end{tabular}
\caption{Croisement complet sur $h = 7$ (126 graphes).}
\end{table}

\section{Discussion}

\paragraph{La symetrie 5/7 n'apporte rien dans ce regime.}
Les configurations \texttt{baseline\_v2} et \texttt{sym1} donnent
exactement les memes resultats sur $h = 7$ (29\,663 sols, 50.5~\%). Le
filtrage CSP en amont (notamment via les flags no-freeze + no-table)
produit deja des compositions globalement equilibrees en pentagones et
heptagones ; le cap C-SYM est donc redondant. La loi empirique observee
sur \texttt{db\_v2} reflete une correlation, pas une marge de
manoeuvre exploitable a ce stade du pipeline.

\paragraph{La contrainte pent-en-bord est le levier principal.}
\texttt{pb1} domine systematiquement sur tous les $h$, avec un gain de
+8 a +17 points de taux de planarite par rapport a \texttt{pb2}. Le
choix $K_{\text{pb}} = 1$ s'avere optimal : \texttt{pb2} (laxiste) ne
filtre pas assez, \texttt{pb0} (strict) rend le probleme insatisfiable
sur tous les $h$ testes.

\paragraph{Trade-off volume / qualite.}
\texttt{pb1} divise environ par 7 le nombre absolu de solutions
generees par rapport a \texttt{pb2}, mais multiplie le taux de
planarite par 1.3--1.5. Pour un usage aval avec validation manuelle, le
ratio signal/bruit de \texttt{pb1} est nettement preferable. Le volume
total de plans pb1 sur $h \in \{6,...,9\}$ reste de l'ordre de
53\,000 molecules, largement suffisant pour la suite.

\paragraph{Decroissance avec $h$.}
Le taux pb1 passe de 80.5~\% ($h{=}6$) a 31.3~\% ($h{=}9$). Au-dela de
l'effet combinatoire bord/interieur, des facteurs purement geometriques
(courbure cumulee, frustration sterique) prennent le pas pour les
grosses molecules. La contrainte structurelle seule ne suffit plus a
elle seule a ces tailles ; une iteration future pourrait combiner C-PB
avec un critere geometrique a priori (par exemple sur la distribution
des angles dieder).

\section{Conclusion}

L'experimentation valide quantitativement la deuxieme loi empirique
identifiee par le data mining : le nombre de pentagones en bord est un
predicteur causal du taux de planarite, et le serrer en amont via une
contrainte CSP ($K_{\text{pb}} = 1$) ameliore directement la qualite
des molecules generees, sans degrader le volume au point de le rendre
inexploitable. La premiere loi (symetrie 5/7) s'avere en revanche non
actionnable a ce stade du pipeline : elle est deja satisfaite par les
contraintes existantes.

Le pipeline DB-only complet (worker DB SQLite + merge idempotent)
fonctionne sans erreur sur les 16 noeuds du cluster (0 echec sur
8\,514 jobs en production), et constitue une base reutilisable pour
les iterations suivantes.

\end{document}
